\documentclass[11pt]{article}
\usepackage{ochamath}
\subjectcolor{2F5DA8}
\setsheet{線形代数・電気系院試補充}{文字パラメータと階数　演習}{https://university-math-crj.pages.dev/linear-algebra/parameter-rank/}
\namelinetrue
\begin{document}
\makesheethead{問題}
自作問題。本文の例題とは数値や設定が異なります。条件・途中式・検算を示してください。
\problem[2つの例外値]
$A(t)=\begin{pmatrix}1&1&0\\1&t&0\\0&0&t+2\end{pmatrix}$、$b=(2,4,0)^{\mathsf T}$。階数と全解を $t\in\mathbb R$ で分類せよ。
\workspace{45mm}
\problem[階数1への落下]
$A(t)=\begin{pmatrix}2&1&1\\2&t&1\\2&1&t\end{pmatrix}$ と任意の $b$ について $t=1$ の可解条件と全解を求めよ。
\workspace{45mm}
\newpage
\problem[左零空間]
$A=\begin{pmatrix}1&0\\0&1\\1&2\end{pmatrix}$ について $Ax=b$ の可解条件を左零空間から求めよ。
\workspace{45mm}
\problem[積の階数下限]
$A\in\mathbb F^{m\times n},B\in\mathbb F^{n\times p}$ に対し $\operatorname{rank}(AB)\ge\operatorname{rank}A+\operatorname{rank}B-n$ を証明せよ。
\workspace{45mm}
\newpage
\problem[上限と下限の達成]
$A=\operatorname{diag}(1,1,0)$。$B_1=\operatorname{diag}(1,1,0)$、$B_2=\operatorname{diag}(0,1,1)$ の積の階数を求めよ。
\workspace{45mm}
\problem[正則変換]
正則な $P,Q$ に対して $\operatorname{rank}(PAQ)=\operatorname{rank}A$ を証明し、$\ker(PAQ)$ を $\ker A$ で表せ。
\workspace{45mm}
\end{document}
